\documentclass[12pt,a4paper]{article}

% ============ Packages ============
\usepackage[UTF8]{ctex}
\usepackage[top=2.5cm, bottom=2.5cm, left=2.5cm, right=2.5cm]{geometry}
\usepackage{amsmath,amssymb,amsthm}
\usepackage{graphicx}
\usepackage{booktabs}
\usepackage{caption}
%\usepackage{multirow}
\usepackage{hyperref}
\usepackage{algorithm}
\usepackage{algpseudocode}
\usepackage{xcolor}
\usepackage{enumerate}
\usepackage{float}
\usepackage{cite}
\usepackage{bm}
\usepackage{listings}

\lstset{
  basicstyle=\footnotesize\ttfamily,
  numbers=left,
  numberstyle=\tiny\color{gray},
  numbersep=8pt,
  frame=single,
  breaklines=true,
  tabsize=4,
  showstringspaces=false,
  keywordstyle=\color{blue}\bfseries,
  commentstyle=\color{green!40!black},
  stringstyle=\color{orange},
  captionpos=b,
  language=Python,
}

\floatname{algorithm}{算法}
\makeatletter
\newenvironment{breakablealgorithm}
{%
  \par\noindent
  \refstepcounter{algorithm}%
  \hrule height .8pt depth 0pt \kern 2pt
  \renewcommand{\caption}[2][\relax]{%
    {\raggedright\textbf{\fname@algorithm~\thealgorithm}\quad ##2\par}%
    \ifx\relax##1\relax
    \addcontentsline{loa}{algorithm}{\protect\numberline{\thealgorithm}##2}%
    \else
    \addcontentsline{loa}{algorithm}{\protect\numberline{\thealgorithm}##1}%
    \fi
    \kern 2pt\hrule\kern 2pt
  }%
}
{\kern 2pt\hrule\par}
\makeatother

% ============ Theorem environments ============
\newtheorem{theorem}{定理}[section]
\newtheorem{lemma}[theorem]{引理}
\newtheorem{definition}[theorem]{定义}
\newtheorem{proposition}[theorem]{命题}
\newtheorem{corollary}[theorem]{推论}

% ============ Title ============
\title{\textbf{浙江省市县联赛分组、抽签、选址与赛制优化模型}}
\date{}

\begin{document}

\maketitle

\begin{abstract}
  浙江省市县联赛拟由11个设区市、20个县级市、33个县共64个单位参赛，分为16个小组进行单循环比赛后决出32强淘汰晋级。赛事组织面临分组公平性、抽签随机性、选址合理性和赛制竞技保真度四类核心问题。本文围绕上述问题建立了从分组、抽签、选址到赛制评价的完整数学模型体系，主要工作如下。

  \textbf{分组方案}：将分组建模为约束满足问题，定义C3冲突对数$F_1$、实力标准差$F_2$和实力极差$F_2'$三项指标，提出贪心优先填充、蛇形轮转、ILP-C3最小化和ILP均衡优化四种策略。利用Hall婚配定理严格证明满足C1、C2硬约束的分组方案恒存在，并由ILP均衡方案实现$F_1=0$、$F_2=0.484$、$F_2'=1$的最优均衡，各组实力极差仅1个实力分。

  \textbf{抽签机制}：提出前瞻最小费用流辅助抽签法，核心思想是在每步抽签前通过求解剩余最小费用流，仅保留能延续全局最优分配的小组进入抽签盆，从源头上杜绝死锁。分析了传统贪心抽签的死锁机制与触发条件（1\,000种子实验中死锁率9.9\%），前瞻流方案则实现零失败和$F_1=0$。在100个随机行政拓扑的广义检验中，前瞻网络流在全部300次抽签中失败率为零（贪心死锁率13\%），当C3零冲突在数学上不可实现时自动转为最小冲突分配。

  \textbf{比赛地点}：将选址建模为带主场禁止约束的设施选址与容量指派ILP模型。基准方案（haversine距离）得到78\,01\,km总旅行距离的最优选址；进一步构造公路时间、铁路时间、综合时间、人口加权和影响力Top20候选集五种成本模型检验选址稳健性，推荐以综合时间模型为主要选址依据。

  \textbf{赛制优化}：基于Bradley-Terry模型和层内正态随机实力赋值进行5\,000次赛事仿真。结果表明，将单循环改为双循环可使Spearman相关系数由0.471提升至0.564，Top-32晋级准确率由69.2\%提升至74.6\%。结合仿真诊断，提出种子分档、同组回避交叉对阵、小组第三附加赛和双败淘汰试点等增强方案及分阶段实施路径。

  \noindent\textbf{关键词：}约束满足问题；前瞻最小费用流；Hall婚配定理；设施选址；Bradley-Terry模型；赛制仿真
\end{abstract}

\newpage

\section{问题重述}

\subsection{问题背景}

2025年江苏城市足球联赛（“苏超”）引发了全社会的广泛关注。拟议中的浙江省市县联赛有浙江省11个副省级城市和地级市（合称市级），20个县级市和33个县（合称县级），共64个单位参加。比赛分为两个阶段：第一阶段为小组赛，参赛队分为16个小组，每个小组4支队伍，进行单循环比赛，小组前两名进入下一阶段；第二阶段为淘汰赛，32支队伍经过5轮比赛，决出冠军。

\subsection{分组约束}

体育赛事分组应遵循公平、透明的原则，本次赛事分组还需遵循以下原则：

{\renewcommand{\labelenumi}{(\textbf{C\arabic{enumi}})}
  \begin{enumerate}
    \item 各市级队应分配在不同小组；
    \item 市级队不与该市代管县的县级队分配在同一小组；
    \item 同一市代管县的县级队应尽量不分配在同一小组。
\end{enumerate}}

\subsection{行政区划数据}

浙江省11个设区市及其代管县市的数量分布如表\ref{tab:city_data}所示。

\begin{table}[H]
  \centering
  \caption{浙江省各市球队分布}
  \label{tab:city_data}
  \small
  \setlength{\tabcolsep}{3pt}
  \begin{tabular}{lcccc}
    \toprule
    设区市 & 市级队 & 代管县级市 & 代管县 & 县级队合计 \\
    \midrule
    杭州 & 杭州市 & 建德市 & 桐庐县、淳安县 & 3 \\
    宁波 & 宁波市 & 余姚市、慈溪市 & 象山县、宁海县 & 4 \\
    温州 & 温州市 & 瑞安市、乐清市、龙港市 & 永嘉县、平阳县等5县 & 8 \\
    嘉兴 & 嘉兴市 & 海宁市、平湖市、桐乡市 & 嘉善县、海盐县 & 5 \\
    湖州 & 湖州市 & — & 德清县、长兴县、安吉县 & 3 \\
    绍兴 & 绍兴市 & 诸暨市、嵊州市 & 新昌县 & 3 \\
    金华 & 金华市 & 兰溪市、义乌市、东阳市、永康市 & 武义县、浦江县、磐安县 & 7 \\
    衢州 & 衢州市 & 江山市 & 常山县、开化县、龙游县 & 4 \\
    舟山 & 舟山市 & — & 岱山县、嵊泗县 & 2 \\
    台州 & 台州市 & 温岭市、临海市、玉环市 & 三门县、天台县、仙居县 & 6 \\
    丽水 & 丽水市 & 龙泉市 & 青田县、缙云县等7县 & 8 \\
    \midrule
    合计 & 11 & 20 & 33 & 53 \\
    \bottomrule
  \end{tabular}
\end{table}

\subsection{待解决的问题}

{\renewcommand{\labelenumi}{\textbf{问题\arabic{enumi}}}
  \begin{enumerate}
    \item 给出若干符合要求的分组方案，并从不同维度比较它们的优劣。
    \item 设计一种可行的抽签方案，使抽签得到的分组方案满足相应要求，并具有较好的性质。
    \item 选择八个市或县作为小组赛的比赛地点（每个地点进行两个小组的比赛），考虑公平性、赛事影响和普及面、交通便利性。
    \item 基于数学模型和定量分析，给出对赛制的合理化建议。
\end{enumerate}}

\section{基本假设与符号说明}

\subsection{基本假设}

{\renewcommand{\labelenumi}{\textbf{(\arabic{enumi})}}
  \begin{enumerate}
    \item 题目未给出各队真实竞技水平，采用层级赋权法近似实力：市级队$s=3$，县级市队$s=2$，县队$s=1$。该赋值反映行政层级与资源禀赋的一般差异，作为全文公平性和均衡分析的基准。后文分别通过问题2的随机拓扑实验和问题4的层内正态随机实力模型（式(\ref{eq:random_strength})）检验核心结论的稳健性。
    \item 基准选址模型采用球面坐标下的大圆距离（haversine公式）计算旅行距离；扩展模型进一步引入公路、铁路旅行时间矩阵。
    \item 城市影响力和赛事普及面用常住人口、交通可达性等公开指标近似刻画。
    \item 单场胜负结果服从Bradley-Terry模型，强制分出胜负（无平局）。真实足球比赛存在平局和积分3/1/0制，本模型将其简化为必分胜负的BT概率，会高估单场比赛的信息量，因此本文只比较同一简化模型下不同赛制的相对差异。小组同分排序中的进球数用Poisson随机量近似。
\end{enumerate}}

\subsection{符号说明}

\begin{table}[H]
  \centering
  \caption{主要符号说明}
  \begin{tabular}{cl}
    \toprule
    符号 & 含义 \\
    \midrule
    $x_{ij}$ & 球队$i$是否分入小组$j$的0-1变量 \\
    $n_{cj}$ & 城市$c$在小组$j$中的县级队数量 \\
    $s_i$ & 球队$i$的实力分值（市=3，县级市=2，县=1） \\
    $S_j$ & 小组$j$的实力总和 \\
    $F_1$ & C3冲突对数（同市县级队同组的无序对数） \\
    $F_2$ & 16组实力标准差 \\
    $F_2'$ & 16组实力极差（最大减最小） \\
    $\mathcal{G}_j$ & 第$j$个小组的球队集合 \\
    $y_k$ & 地点$k$是否被选为比赛场地的0-1变量 \\
    $z_{gk}$ & 小组$g$是否安排到场地$k$的0-1变量 \\
    $d_{ik}$ & 球队$i$所在地到场地$k$的距离 \\
    $c_{ik}$ & 球队$i$到场地$k$的等效成本 \\
    $\tau$ & 约束紧度 = $\max(n_i)/15$ \\
    \bottomrule
  \end{tabular}
\end{table}

\section{问题1：分组方案设计与比较}

\subsection{数学模型}

\subsubsection{决策变量与基本约束}

定义0-1决策变量$x_{ij} \in \{0,1\}$，表示球队$i$是否分入小组$j$（$i = 1,2,\dots,64$，$j = 1,2,\dots,16$）。

每支球队恰好进入一个小组，每个小组恰好容纳4支球队：

\begin{equation}
  \sum_{j=1}^{16} x_{ij} = 1,\quad \forall i; \qquad \sum_{i=1}^{64} x_{ij} = 4,\quad \forall j
  \label{eq:assign}
\end{equation}

\textbf{C1（市级队不同组）}：令$M$为11支市级队的集合，则

\begin{equation}
  \sum_{i \in M} x_{ij} \leq 1,\quad \forall j
  \label{eq:C1}
\end{equation}

\textbf{C2（市级队不与本市县级队同组）}：令$m(c)$为城市$c$的市级队，$C_c$为城市$c$下辖县级队集合。则

\begin{equation}
  x_{m(c),j} + x_{i,j} \leq 1,\quad \forall i \in C_c,\ \forall c \in \{1,\dots,11\},\ \forall j
  \label{eq:C2}
\end{equation}

\textbf{C3（同市县级队尽量不同组）}：C3是软约束。若将其严格化，则

\begin{equation}
  \sum_{i \in C_c} x_{ij} \leq 1,\quad \forall c,\ \forall j
  \label{eq:C3_hard}
\end{equation}

\subsubsection{容量与可行域分析}

16组总容量为$16 \times 4 = 64$，恰好容纳全部球队。其中11组含有市级队（各剩3个县级名额），5组无市级队（各4个县级名额），县级名额总计$11 \times 3 + 5 \times 4 = 53$，恰好等于县级队总数。

\textbf{C3可行性}：各市县级队需尽量分散到不同组。关键约束来自县级队最多的两个城市——温州（8支）和丽水（8支）。每市仅需避开本市市级队所在的1个组，可用组数为15。由于$8 \leq 15$，单个城市层面不存在组数瓶颈。下文先用Hall婚配定理证明C1+C2硬约束恒可行，再用ILP求解结果给出本题C3零冲突的构造性证据。

\subsubsection{基于Hall婚配定理的严格证明}

上述直观论证可以由Hall婚配定理给出严格的形式化证明。

\begin{definition}[约束二部图]
  构造二部图$G = (U \cup V, E)$，其中：
  \begin{itemize}
    \item $U$：53支县级队的集合；
    \item $V$：53个县级队槽位（有市级队的组各3个，无市级队的组各4个）；
    \item $(u, v) \in E$当且仅当：县级队$u$所属市的市级队不在槽位$v$所在的组中（即不违反C2约束）。
  \end{itemize}
\end{definition}

在该二部图中，完美匹配对应一个满足C1和C2的完整分组方案（注意：该二部图仅编码C2禁入关系，不含C3同市互斥约束；C3的可满足性由下文推论给出论证）。

\begin{proposition}[C1+C2可行性]
  固定任意一种满足C1的市级队落位。对任意县级队子集$S\subseteq U$，记$B(S)$为$S$中球队所属城市的集合。若$|B(S)|=1$，则$S$中所有队伍只共同禁入1个含市级队的小组，邻域至少包含其余15个小组的50个县级槽位，因此$|N(S)|=50\ge |S|$。若$|B(S)|\ge 2$，则任意一个县级槽位至多只会被其所属城市的县级队禁入，但一定可由$S$中其他城市的县级队使用，所以$N(S)=V$，有$|N(S)|=53\ge |S|$。于是Hall条件对所有$S\subseteq U$成立，二部图$G$存在完美匹配，即满足C1和C2的完整分组方案一定存在。
\end{proposition}

\begin{corollary}[本题C3零冲突的构造性证据]
  C3是软约束，Hall证明本身不包含“同市县级队互斥”条件。本文将C3冲突对数$F_1$纳入ILP目标，并用CBC求解得到$F_1=0$的方案D；附录A列出完整分组。该结果构成了本题数据下C3可完全满足的构造性证据。若换成其他行政拓扑，C3零冲突未必可行，因此问题2中将C3作为最小费用目标而不是硬约束处理。
\end{corollary}

Hall定理保证了全局解的存在性，但不保证任意贪心决策路径都能通向可行解。问题2中将详细分析贪心抽签算法在何种条件下会陷入死锁，并在此基础上设计一种“事前筛选”的前瞻抽签机制。

\subsection{评价指标体系}

\textbf{指标1：C3冲突对数$F_1$}

\begin{equation}
  F_1 = \sum_{c=1}^{11} \sum_{j=1}^{16} \binom{n_{cj}}{2}
  \label{eq:F1}
\end{equation}

其中$n_{cj} = \sum_{i \in C_c} x_{ij}$为城市$c$分入小组$j$的县级队数量。$F_1 = 0$表示C3完全满足，越小越好。

\textbf{指标2：实力标准差$F_2$与极差$F_2'$}

各组实力总和$S_j = \sum_{i} s_i x_{ij}$（市级$s=3$，县级市$s=2$，县$s=1$），总实力$\sum s_i = 106$，组均值$\bar{S} = 6.625$。

\begin{equation}
  F_2 = \sqrt{\frac{1}{16} \sum_{j=1}^{16} (S_j - \bar{S})^2}, \qquad
  F_2' = \max_j S_j - \min_j S_j
  \label{eq:F2}
\end{equation}

两者均越小越好，反映各组实力是否均衡。

\subsection{分组方案生成算法}

\subsubsection{方案A：贪心优先填充}

\textbf{思想}：按县级队数量从多到少依次处理各市，每支县级队贪心选择C3冲突最小且人数最少的组。

\begin{breakablealgorithm}
  \caption{方案A：贪心优先填充}
  \label{alg:q1_greedy}
  \begin{algorithmic}[1]
    \State \textbf{输入：} 11个城市的市级队和县级队，16个小组，组容量4
    \State \textbf{输出：} 分组方案$\{\mathcal{G}_1,\dots,\mathcal{G}_{16}\}$
    \State 初始化$\mathcal{G}_1,\dots,\mathcal{G}_{16}\gets\emptyset$
    \State 按县级队数降序排列城市，得到序列$C$
    \For{$r=1$ \textbf{to} $11$}
    \State 将城市$C_r$的市级队放入$\mathcal{G}_r$，记录禁入组$m(C_r)=r$
    \EndFor
    \For{城市$c\in C$}
    \State 随机打乱城市$c$的县级队序列$T_c$
    \For{县级队$t\in T_c$}
    \State $A\gets\{j:\ |\mathcal{G}_j|<4,\ j\neq m(c)\}$
    \State 对每个$j\in A$计算$\mathrm{same}_{cj}$：组$j$中是否已有城市$c$的县级队
    \State 选择$j^*=\arg\min_{j\in A}(\mathrm{same}_{cj},|\mathcal{G}_j|)$；若并列则随机选取
    \State 将$t$放入$\mathcal{G}_{j^*}$
    \EndFor
    \EndFor
    \State \Return $\{\mathcal{G}_j\}_{j=1}^{16}$
  \end{algorithmic}
\end{breakablealgorithm}

\subsubsection{方案B：蛇形轮转分配}

\textbf{思想}：借鉴体育赛事中的“蛇形排列”，将各市县级队在小组间按轮转方式分散放置。

\begin{breakablealgorithm}
  \caption{方案B：蛇形轮转分配}
  \label{alg:q1_serpentine}
  \begin{algorithmic}[1]
    \State \textbf{输入：} 城市序列$C$，16个小组，组容量4
    \State \textbf{输出：} 分组方案$\{\mathcal{G}_1,\dots,\mathcal{G}_{16}\}$
    \State 初始化小组并按算法\ref{alg:q1_greedy}第3--6行放置市级队
    \State 构造蛇形序列$P=(1,2,\dots,16,16,15,\dots,1,\dots)$，令指针$p\gets1$
    \For{城市$c\in C$}
    \State 随机打乱城市$c$的县级队序列$T_c$，并置$U_c\gets\emptyset$
    \For{县级队$t\in T_c$}
    \State $\mathrm{placed}\gets\mathrm{false}$
    \For{$\ell=1$ \textbf{to} $16$}
    \State $j\gets P_p$，$p\gets p+1$
    \If{$j\neq m(c)$ \textbf{and} $|\mathcal{G}_j|<4$ \textbf{and} $j\notin U_c$}
    \State 将$t$放入$\mathcal{G}_j$，$U_c\gets U_c\cup\{j\}$，$\mathrm{placed}\gets\mathrm{true}$
    \State \textbf{break}
    \EndIf
    \EndFor
    \If{\textbf{not} $\mathrm{placed}$}
    \State 在满足C2且未满的组中顺延放置$t$；优先选择$j\notin U_c$的组
    \EndIf
    \EndFor
    \EndFor
    \State \Return $\{\mathcal{G}_j\}_{j=1}^{16}$
  \end{algorithmic}
\end{breakablealgorithm}

\subsubsection{方案C：ILP-C3最小化}

以$F_1$最小化为目标，C1和C2为硬约束，建立整数线性规划模型（ILP）：

\begin{align}
  \min \quad & \sum_{c=1}^{11} \sum_{j=1}^{16} \binom{n_{cj}}{2} \\
  \text{s.t.} \quad & \text{式}(\ref{eq:assign}), (\ref{eq:C1}), (\ref{eq:C2}) \nonumber
\end{align}

代码实现中使用PuLP调用COIN-OR CBC求解器求解。

\subsubsection{方案D：ILP均衡优化}

实际求解时先将C3冲突对数作为高权重目标，再最小化组间实力极差。令同市县级队配对变量$u_{pj}$表示配对$p$是否同落入组$j$，并引入$S_{\max}$、$S_{\min}$约束各组实力和，则目标为

\begin{align}
  \min \quad & w_1\sum_{p,j} u_{pj} + w_2(S_{\max}-S_{\min}) \\
  \text{s.t.} \quad & \text{式}(\ref{eq:assign}), (\ref{eq:C1}), (\ref{eq:C2}) \nonumber \\
  & u_{pj} \ge x_{a_p,j}+x_{b_p,j}-1,\quad \forall p=(a_p,b_p),j \nonumber \\
  & S_{\max}\ge S_j,\quad S_{\min}\le S_j,\quad \forall j . \nonumber
\end{align}

这里我们取$w_1=10,w_2=1$，求解结果自然达到$F_1=0$，同时将实力极差压缩到1。

\subsection{实验结果与方案对比}

四种方案的评价指标对比如表\ref{tab:q1_comparison}所示。附录A列出方案A--D的完整分组，便于直接核查各方案的约束满足情况和实力分布。

\begin{table}[H]
  \centering
  \caption{四种分组方案的评价指标对比}
  \label{tab:q1_comparison}
  \small
  \begin{tabular}{lccc}
    \toprule
    方案 & $F_1$（C3冲突对数） & $F_2$（实力$\sigma$） & $F_2'$（极差） \\
    \midrule
    方案A：贪心优先填充 & 0 & 0.927 & 3 \\
    方案B：蛇形轮转分配 & 0 & 1.111 & 4 \\
    方案C：ILP-C3最小化 & 0 & 0.992 & 3 \\
    方案D：ILP均衡优化 & 0 & \textbf{0.484} & \textbf{1} \\
    \bottomrule
  \end{tabular}
\end{table}

\textbf{分析}：

\begin{enumerate}[(1)]
  \item \textbf{C3满足度}：四种方案均实现了$F_1=0$，验证了“C3在本题数据下可完全满足”的预判。
  \item \textbf{实力均衡度}：ILP均衡方案（方案D）表现最优，$F_2=0.484$仅为贪心方案（0.927）的52\%，$F_2'=1$意味着各组实力极差仅为1个实力分，实现了高度的竞技公平。
  \item \textbf{方案推荐}：若优先追求公平性，应选方案D（ILP均衡）；若兼顾可解释性和计算效率，方案A（贪心）是合理的次优选择。
\end{enumerate}

\section{问题2：抽签方案设计}

\subsection{设计目标}

抽签方案需满足以下性质：
\begin{enumerate}[(1)]
  \item \textbf{可行性}：抽签结果必须满足C1、C2硬约束，并尽量满足C3。
  \item \textbf{随机性}：在约束允许范围内保留尽可能多的随机选择；前瞻流方案先随机抽队，再从安全小组集合中随机抽取小组。
  \item \textbf{透明性}：抽签流程公开，每一步可审计。
  \item \textbf{鲁棒性}：不会因随机因素陷入死锁（无解）。
\end{enumerate}

\subsection{贪心死锁与前瞻抽签思想}

\subsubsection{朴素贪心抽签}

在介绍完整抽签方案之前，先分析一种朴素贪心策略的死锁问题。该策略与问题1方案A类似，具体流程如算法\ref{alg:q2_naive}所示。

\begin{breakablealgorithm}
  \caption{朴素贪心抽签}
  \label{alg:q2_naive}
  \begin{algorithmic}[1]
    \State \textbf{输入：} 球队数据，随机种子
    \State \textbf{输出：} 分组方案或失败标记$\bot$
    \State 随机将11支市级队分配到16组中的11个不同组，并记录各城市禁入组$m(c)$
    \State 按县级队数量从多到少排列城市，得到序列$C$
    \For{城市$c\in C$}
    \State 随机打乱城市$c$的县级队序列$T_c$
    \For{县级队$t\in T_c$}
    \State $A\gets\{j:\ |\mathcal{G}_j|<4,\ j\neq m(c)\}$
    \If{$A=\emptyset$}
    \State \Return $\bot$
    \EndIf
    \State 对每个$j\in A$计算C3冲突标记$\mathrm{same}_{cj}$
    \State 选择$j^*=\arg\min_{j\in A}(\mathrm{same}_{cj},|\mathcal{G}_j|)$；若并列则随机选取
    \State 将$t$放入$\mathcal{G}_{j^*}$
    \EndFor
    \EndFor
    \State \Return $\{\mathcal{G}_j\}_{j=1}^{16}$
  \end{algorithmic}
\end{breakablealgorithm}

该策略不包含任何回溯机制，一旦某支队伍无合法候选组，算法即宣告失败。

\subsubsection{全局可行性与贪心死锁的区别}

问题1中已由Hall婚配定理证明：满足C1和C2的分组方案恒存在。然而，贪心策略每一步只看当前的局部最优，一旦放入便不再调整。这种不可逆的局部决策可能导致后续队伍的合法位置被提前占满，使算法陷入"死锁"。

\subsubsection{死锁反例}

以$\text{seed}=0$为例，可以稳定复现一个死锁场景。市级队随机分配结果如表\ref{tab:deadlock_mun}所示。

\begin{table}[H]
  \centering
  \caption{死锁反例中的市级队分配（seed=0）}
  \label{tab:deadlock_mun}
  \begin{tabular}{cccc}
    \toprule
    组号 & 市级队 & 组号 & 市级队 \\
    \midrule
    G1 & 金华市 & G7 & 绍兴市 \\
    G2 & 丽水市 & G8 & 温州市 \\
    G3 & 杭州市 & G9 & 宁波市 \\
    G4 & 湖州市 & G10 & 台州市 \\
    G5 & 衢州市 & G11 & \textbf{舟山市} \\
    G6 & 嘉兴市 & G12--G16 & 无市级队 \\
    \bottomrule
  \end{tabular}
\end{table}

算法按县级队数量降序依次处理温州(8)→丽水(8)→金华(7)→台州(6)→嘉兴(5)→衢州(4)→宁波(4)→湖州(3)→绍兴(3)→杭州(3)，共放置52支县级队。舟山市仅2支县级队，排在最后。处理过程中，舟山的岱山县已被分配到G7，而G11（舟山市所在组）仅放入3队（舟山市+其他市的县级队），恰好留有1个空位。

当轮到舟山的最后1支队伍嵊泗县时：

\begin{itemize}
  \item G11（舟山市所在组）：C2禁入，不可放入；
  \item G1--G10、G12--G16：全部恰好满员（4队）。
\end{itemize}

唯一有空位的组是禁入组，全部合法位置均被占用——算法陷入不可恢复的死锁。这是因为贪心策略在处理前10个市的52支县级队时均匀填充了所有组，不预判后续城市的需求，未为最后处理的舟山市保留必要的灵活性。

\subsubsection{死锁的触发条件与频率}

1\,000个固定随机种子实验中，朴素贪心策略死锁率为9.90\%。死锁更易在以下条件下触发：

\begin{enumerate}[(1)]
  \item 县级队数量少的城市被排在最后处理；
  \item 前续城市的贪心放置恰好填满了所有非禁入组；
  \item 最后待分配队伍的C2禁入组是唯一有空位的组。
\end{enumerate}

广义蒙特卡洛实验表明，死锁并非由约束系统本身无解导致，而是由抽签路径的局部不可逆性导致：在100个随机拓扑、每个拓扑3个配对种子的实验中，朴素贪心总死锁率为13.0\%，而前瞻网络流失败率为0（见表\ref{tab:generalized}）。

\subsection{前瞻最小费用流辅助抽签法}

朴素贪心只依据当前局部信息决策，遇到死锁后只能整体作废重抽。若用于真实体育抽签仪式，已经公布的结果不宜被撤回，也不宜在事后调整已公布分组。因此采用事前筛选思想：每抽出一支县级队后，计算机先判断哪些小组能够通向全局最优的剩余分配，再把这些安全小组编号放入现场小组盆中进行物理抽取。

\subsubsection{状态变量}

设当前抽签状态下：
\begin{itemize}
  \item $r_j$：小组$j$的剩余名额；
  \item $m(c)$：城市$c$的市级队所在小组，即城市$c$县级队的C2禁入组；
  \item $a_{cj}$：城市$c$已有多少支县级队在小组$j$中；
  \item $b_c$：城市$c$尚未抽出的县级队数量。
\end{itemize}

当前C3冲突对数为
\begin{equation}
  F_1^{\mathrm{cur}} = \sum_c\sum_{j=1}^{16}\binom{a_{cj}}{2}.
\end{equation}
若未来还要把$y_{cj}$支城市$c$的县级队放入小组$j$，则新增C3冲突为
\begin{equation}
  \Delta_{cj}(y_{cj})
  =
  \binom{a_{cj}+y_{cj}}{2}-\binom{a_{cj}}{2}.
  \label{eq:q2_delta_c3}
\end{equation}
因此，C3不是普通的边可行性约束，而是城市--小组层面的凸费用。

\subsubsection{剩余分配的最小费用流模型}

对当前抽出的县级队$t$（所属城市为$c_t$），假设尝试将其放入小组$g$。若$g=m(c_t)$或$r_g=0$，则该组直接排除；否则临时放入$t$，更新$r_g,a_{c_tg},b_{c_t}$，并对剩余队伍求解如下最小费用流：

\begin{align}
  \min \quad & \sum_c\sum_{j=1}^{16}
  \left[
    \binom{a_{cj}+y_{cj}}{2}-\binom{a_{cj}}{2}
  \right] \label{eq:q2_mcf_obj}\\
  \text{s.t.}\quad
  & \sum_{j=1}^{16} y_{cj}=b_c,\quad \forall c, \label{eq:q2_mcf_city}\\
  & \sum_c y_{cj}=r_j,\quad \forall j, \label{eq:q2_mcf_group}\\
  & y_{cj}=0,\quad j=m(c), \label{eq:q2_mcf_c2}\\
  & y_{cj}\in\mathbb{Z}_{\ge0}. \nonumber
\end{align}

网络构造如下。源点连向每个城市节点$c$，容量$b_c$、费用0；城市$c$连向小组满足 $j \neq m(c)$ 的 $j$的边用单位容量平行边线性化式(\ref{eq:q2_delta_c3})，第$\ell$条边的费用为$a_{cj}+\ell-1$；小组$j$连向汇点，容量$r_j$、费用0。若最大流量等于$\sum_c b_c$，则剩余抽签可完成；其最小费用即为未来最少还需增加的C3冲突数。

为提高实现效率，程序先用最大流检查$F_1=0$是否仍可行：只保留$a_{cj}=0$的城市--小组边，容量为1；若该最大流满流，则该候选组的最终$F_1$可保持为0，无需调用完整最小费用流。只有零冲突不可延续时，才退回到完整的最小费用流模型。

\subsubsection{现场抽签流程}

\begin{breakablealgorithm}
  \caption{前瞻最小费用流辅助抽签法}
  \label{alg:q2_lookahead_flow}
  \begin{algorithmic}[1]
    \State \textbf{输入：} 11支市级队、53支县级队、16个小组、组容量4
    \State \textbf{输出：} 完整分组方案
    \State \textbf{阶段1：市级队定位}
    \State 从市级队盆中随机抽取11支市级队，并随机抽取11个不同小组放入，记录$m(c)$
    \State \textbf{阶段2：县级队前瞻抽签}
    \State 将53支县级队混合放入县级队盆，随机抽出一支队$t$，记其城市为$c_t$
    \For{每个小组$g=1,\dots,16$}
    \If{$g=m(c_t)$ \textbf{or} $r_g=0$}
    \State 标记$g$为不可用
    \Else
    \State 临时将$t$放入$g$，求解剩余队伍的最大流/最小费用流
    \If{剩余最大流不能满流}
    \State 标记$g$为不可完成
    \Else
    \State 记录该选择对应的最小最终C3冲突数$\Phi(g)$
    \EndIf
    \State 撤销临时放置
    \EndIf
    \EndFor
    \State 令$\Phi^*=\min_g\Phi(g)$，安全小组集$A(t)=\{g:\Phi(g)=\Phi^*\}$
    \State 将$A(t)$中的小组编号放入现场小组盆，随机抽出正式小组
    \State 重复阶段2，直到53支县级队全部抽取完毕
    \State \Return 16个小组
  \end{algorithmic}
\end{breakablealgorithm}

这里的最小费用最大流在我们的代码中使用 Bellman-Ford 和 Dinic 算法实现。

\subsubsection{可行性保证}

\begin{proposition}[前瞻抽签不死锁]
  若当前部分抽签状态存在至少一个满足C1、C2的完整延拓，则算法\ref{alg:q2_lookahead_flow}在当前抽出的球队$t$上得到的安全小组集$A(t)$非空。若从$A(t)$中任取一个小组正式放置$t$，则新状态仍存在满足C1、C2的完整延拓；并且该选择在所有可完成选择中保持最终$F_1$的最小可实现值。
\end{proposition}

\textbf{证明}。设当前状态存在一个完整延拓。由于$t$是尚未放置的球队之一，在该延拓中$t$必被放入某个小组$g^\star$。该小组必然满足容量约束且不等于$t$所属城市的C2禁入组；固定选择$t\to g^\star$后，其他未放置球队仍可按照该延拓完成分配。因此，算法枚举小组时，$g^\star$对应的剩余流模型能够满流，$\Phi(g^\star)$有定义，故至少存在一个可完成候选组，$A(t)$非空。

对任意$g\in A(t)$，其定义已经要求固定$t\to g$后的剩余流模型可满流，因此正式选择该组后仍存在完整延拓，不会进入死锁状态。同时，$A(t)$只保留使$\Phi(g)$达到最小值$\Phi^\ast=\min_g\Phi(g)$的候选组，所以该选择在所有可完成选择中保持最终C3冲突数最小。由归纳法，从初始存在完整延拓的状态出发，每一步之后仍存在完整延拓，直到53支县级队全部放置完成，算法均不会出现无可用小组。 \qed

该命题说明，前瞻流不是事后调整或整体重抽，而是在每一步把会导致未来死锁或劣化C3目标的危险小组提前从现场小组盆中剔除。由于C3以费用而非硬禁令进入模型，当某些泛化数据无法实现$F_1=0$时，该方法仍会自动选择C3冲突数最小的可完成路径。

\subsection{策略对比实验}

为验证新抽签法的效果，本文在同一批1\,000个随机种子上比较朴素贪心与前瞻最小费用流。朴素贪心代表“不预判未来”的传统现场抽签，前瞻流代表本文推荐的事前筛选机制。由于前瞻流每步需要求解多个流模型，计算成本高于贪心策略，但本题规模较小，仍能满足现场实时筛选需求。

\begin{table}[H]
  \centering
  \caption{朴素贪心与前瞻流的1\,000个配对种子实验对比}
  \label{tab:strategy_comparison}
  \small
  \begin{tabular}{lcc}
    \toprule
    指标 & 朴素贪心 & 前瞻最小费用流 \\
    \midrule
    成功率 & 90.1\% & \textbf{100\%} \\
    死锁/失败率 & 9.9\% & \textbf{0\%} \\
    $P(F_1=0)$ & 100\%* & \textbf{100\%} \\
    $\mathbb{E}[F_1]$ & 0* & \textbf{0} \\
    $\mathbb{E}[F_2]$ & 1.139 & \textbf{1.129} \\
    $\mathbb{E}[F_2']$ & 4.01 & \textbf{3.98} \\
    平均前瞻检查次数 & -- & 582.7 \\
    平均计算时间/次 & 0.0011s & 0.2153s \\
    \bottomrule
    \multicolumn{3}{l}{\footnotesize *朴素贪心的$F_1,F_2,F_2'$仅统计成功（未死锁）的样本。}
  \end{tabular}
\end{table}

实验结果表明：前瞻最小费用流在样本中实现了零失败和$F_1=0$，同时不需要撤回已经公布的抽签结果。其代价是计算量较高：1\,000次实验中平均每次抽签进行了582.7次前瞻可行性检查，其中58.6次需要调用完整最小费用流模型；单次抽签平均耗时约0.22秒，仍可满足现场“抽队伍--显示可用小组--抽小组”的实时需求。

\subsection{边际均匀性检验}

前瞻流会根据剩余可完成性动态删去部分候选组，因此需要检查该筛选机制是否明显破坏抽签的随机性。本文采用边际均匀性检验：对每支球队$t$，统计其在1\,000次前瞻流抽签中落入16个小组的频数$O_{t1},\dots,O_{t,16}$。在零假设$H_0$下，球队$t$落入各小组的边际概率相等，即期望频数为$E_{tj}=1000/16$。检验统计量为
\begin{equation}
  \chi_t^2=\sum_{j=1}^{16}\frac{(O_{tj}-E_{tj})^2}{E_{tj}},\quad df=15.
\end{equation}

逐队卡方检验中，前瞻流有59/64支球队通过$\alpha=0.05$显著性水平；考虑到同时进行了64次检验，采用Bonferroni校正后显著性阈值为$0.05/64\approx0.0008$，此时64支球队全部通过。进一步将64个$p$值与$U(0,1)$分布做Kolmogorov--Smirnov检验，得到$D=0.0947,p=0.5816$，未发现$p$值整体偏离均匀分布的证据。该结果说明，前瞻筛选虽然会剔除危险小组，但在本题样本规模下未表现出系统性的边际分组偏差。

\subsection{广义蒙特卡洛鲁棒性检验}

前述实验固定于浙江省实际行政区划。为检验抽签死锁机制是否依赖特定数据，本文构造100个随机行政拓扑，并在每个拓扑上运行3个配对随机种子。实验仅比较两类模型：基础贪心抽签与前瞻网络流分配模型。为衡量问题难度，引入约束紧度（tightness）$\tau$：

\begin{equation}
  \tau = \frac{\max_i n_i}{15}
  \label{eq:tau}
\end{equation}

其中$\max_i n_i$为县级队最多的市的县级队数，15为排除C2禁入组后的可用组数。$\tau$的物理意义是：最紧张的城市，其县级队需求占可用组数的比例。若$\tau < 1$，C3有望完全满足；若$\tau = 1$，C3恰好满足但零冗余；若$\tau > 1$，C3零冲突不可能完全满足。浙江省的$\tau = 8/15 = 0.533$，属于中等偏紧水平。

随机拓扑生成时，先设定市级队数量$k\in[5,16]$，再将$64-k$支县级队按不同集中度分配给各市。县级队分布通过Dirichlet分布控制：均匀型采用确定性等分，中等变异采用对称Dirichlet参数$\alpha=(2,\dots,2)$，集中型采用少数城市$\alpha=3$、多数城市$\alpha=0.5$的非均匀参数。每个拓扑上只运行基础贪心和前瞻网络流两个模型。前瞻网络流在$\tau \le 1$时检验能否维持$F_1=0$，在$\tau>1$时则检验其是否能自动转为最小C3冲突分配。

\begin{table}[H]
  \centering
  \caption{不同紧度区间下的基础贪心与前瞻网络流对比（100个随机拓扑）}
  \label{tab:generalized}
  \small
  \setlength{\tabcolsep}{3pt}
  \resizebox{\textwidth}{!}{%
    \begin{tabular}{lcccccc}
      \toprule
      紧度区间 & 拓扑数 & 贪心死锁率 & 贪心$P(F_1{=}0)$* & 流失败率 & 流$P(F_1{=}0)$ & 流$\mathbb{E}[F_1]$ \\
      \midrule
      中等（$\tau {\leq} 0.5$） & 41 & 22.0\% & 100.0\% & 0.0\% & 100.0\% & 0.00 \\
      偏紧（$0.5 {<} \tau {\leq} 0.7$） & 17 & 3.9\% & 100.0\% & 0.0\% & 100.0\% & 0.00 \\
      紧张（$0.7 {<} \tau {\leq} 1.0$） & 16 & 14.6\% & 97.6\% & 0.0\% & 100.0\% & 0.00 \\
      超紧（$\tau {>} 1.0$） & 26 & 3.8\% & 0.0\% & 0.0\% & 0.0\% & 11.50 \\
      \bottomrule
    \end{tabular}
  }
  \vspace{2pt}
  \parbox{\textwidth}{\footnotesize *贪心$P(F_1{=}0)$仅在未死锁的成功样本上统计；当$\tau>1$时，C3零冲突在数学上不可实现。}
\end{table}

表\ref{tab:generalized}说明，基础贪心在各紧度区间均可能死锁，且其死锁率还受到市级队数量、县级队集中度和前序填充路径共同影响，因此$\tau$不应被解释为单一的死锁概率模型。相比之下，前瞻网络流在全部300次抽签中失败率为0；当$\tau\le1$时均保持$F_1=0$，当$\tau>1$导致C3零冲突不可实现时，算法仍能给出可完成的较小冲突方案，平均$F_1=11.50$。
\section{问题3：比赛地点选择}

\subsection{数学模型：设施选址与带容量指派}

题目要求从64个市县中选取8个作为小组赛承办地，每个场地承办2个小组。该问题包含两层决策：选哪些地点开放（选址），以及每个小组去哪个场地（指派）。由于场地集合不同时最优指派也不同，因此将其统一为一个ILP模型同时优化。此外，为避免球队在本地比赛获得主场优势，加入主场禁止约束。

\subsubsection{成本矩阵}

ILP模型只要求给定球队$i$到候选场地$k$的成本矩阵$c_{ik}$。在基准模型中，成本直接取haversine球面距离：

\begin{equation}
  c_{ik}=d_{ik}^{\mathrm{hav}},
  \label{eq:cost}
\end{equation}

其中$d_{ik}^{\mathrm{hav}}$为球队所在地到候选场地的球面距离。主场公平不放入成本函数，而作为硬约束处理：若小组中包含地点$k$的本地球队，则该小组不能安排到地点$k$。后续敏感性分析保持同一ILP结构，仅替换$c_{ik}$为公路时间、铁路时间、综合时间或人口加权时间，并通过Top20候选集体现城市影响力约束。

\subsubsection{ILP模型}

定义决策变量：
\begin{itemize}
  \item $y_k \in \{0,1\}$：地点$k$是否被选为比赛场地（$k = 1,\dots,64$）；
  \item $z_{gk} \in \{0,1\}$：小组$g$是否安排到场地$k$（$g = 1,\dots,16$）。
\end{itemize}

小组$g$若安排到场地$k$，其总成本为$C_{gk} = \sum_{i \in G_g} c_{ik}$。

\begin{align}
  \min \quad & \sum_{g=1}^{16} \sum_{k=1}^{64} z_{gk} C_{gk} \label{eq:q3_obj} \\
  \text{s.t.} \quad & \sum_{k=1}^{64} y_k = 8 \label{eq:q3_y} \\
  & \sum_{k=1}^{64} z_{gk} = 1,\quad \forall g \label{eq:q3_z1} \\
  & z_{gk} \leq y_k,\quad \forall g,k \label{eq:q3_z2} \\
  & \sum_{g=1}^{16} z_{gk} = 2y_k,\quad \forall k \label{eq:q3_z3} \\
  & z_{gk}=0,\quad \text{若地点 } k \in \mathcal{G}_g \label{eq:q3_home}
\end{align}

约束(\ref{eq:q3_y})确保恰好选择8个场地；(\ref{eq:q3_z1})确保每个小组分配到恰好一个场地；(\ref{eq:q3_z2})确保只能分配到已选场地；(\ref{eq:q3_z3})确保每个选定场地恰好承办2个小组；(\ref{eq:q3_home})为主场禁止约束，避免任何球队在本地比赛。

\subsection{实验结果}

由于问题1推荐方案D（ILP均衡优化）为最优分组，选址以方案D的分组为基础进行求解。这里采用 PuLP 库调用 CBC 求解器求解上述ILP模型。求解结果如表\ref{tab:q3_ilp}所示。

\begin{table}[H]
  \centering
  \caption{ILP最优选址结果（基于方案D分组）}
  \label{tab:q3_ilp}
  \begin{tabular}{lcccc}
    \toprule
    场地 & 经度 & 纬度 & 承办小组 & 平均距离/队 \\
    \midrule
    余姚市 & 121.15°E & 30.04°N & 组2、组3 & 158 km \\
    诸暨市 & 120.24°E & 29.71°N & 组5、组6 & 113 km \\
    嵊州市 & 120.83°E & 29.56°N & 组12、组16 & 145 km \\
    东阳市 & 120.24°E & 29.29°N & 组4、组11 & 101 km \\
    武义县 & 119.81°E & 28.90°N & 组8、组10 & 114 km \\
    磐安县 & 120.45°E & 29.06°N & 组1、组15 & 128 km \\
    仙居县 & 120.73°E & 28.85°N & 组13、组14 & 112 km \\
    缙云县 & 120.09°E & 28.66°N & 组7、组9 & 105 km \\
    \bottomrule
  \end{tabular}
\end{table}

选址结果分布在浙中偏南区域（纬28.66°--30.04°N，经119.81°--121.15°E），兼顾了浙北（余姚）、浙中（诸暨、嵊州、东阳、武义、磐安、缙云）和浙南（仙居）。主场禁止约束生效后，最近单队距离由0提升到10 km，避免了本地球队直接获得主场优势。需指出的是，纯距离模型仍偏向省域几何中心，未必自动选择杭州、宁波、温州等核心城市；因此下文进一步引入交通时间、人口加权和影响力候选集。

\subsection{多维成本模型对比}

仅用直线距离会低估海岛、山区和跨水域交通成本，因此进一步构造六类选址实验：haversine距离、公路旅行时间、铁路旅行时间、公铁平均综合时间、人口加权公路时间，以及影响力Top20候选集下的综合时间模型。表\ref{tab:q3_cost_models}给出同一ILP模型在不同成本定义和候选集策略下的求解结果。

\begin{table}[H]
  \centering
  \caption{不同成本模型下的ILP选址结果对比}
  \label{tab:q3_cost_models}
  \small
  \setlength{\tabcolsep}{4pt}
  \begin{tabular}{lcccc}
    \toprule
    成本模型 & 总成本 & 单位 & 平均/队 & 对应haversine距离 \\
    \midrule
    haversine & 7\,801 & km & 121.9 & 7\,801 km \\
    公路时间 & 8\,372 & min & 130.8 & 8\,220 km \\
    铁路时间 & 6\,524 & min & 101.9 & 8\,531 km \\
    综合时间 & 7\,731 & min & 120.8 & 8\,484 km \\
    人口加权 & 662\,682 & 万人$\cdot$min & 10\,354.4 & 9\,119 km \\
    综合时间（影响力Top20） & 8\,130 & min & 127.0 & 8\,990 km \\
    \bottomrule
  \end{tabular}
\end{table}

多维模型显示，不同成本定义会明显改变场地集合：haversine与综合时间模型的场地交集为2/8，而综合时间与影响力Top20综合时间模型的交集为5/8。诸暨市、杭州市、义乌市、绍兴市等地点在多个模型中反复出现，说明这些地点在交通可达性和赛事影响力之间具有较强稳健性。若赛事组织者只追求总旅行距离，可采用haversine或公路时间模型；若更重视真实交通体验，应优先采用综合时间模型；若更重视观众覆盖和传播效果，则可采用人口加权模型或影响力Top20候选集作为约束补充。本文最终建议以综合时间模型为主方案，并将影响力Top20方案作为宣传和承办能力约束下的备选。

\textbf{人口加权成本的定义}：对球队$i$到场地$k$的人口加权成本为$c_{ik}^{\text{pop}} = P_i \cdot t_{ik}$，其中$P_i$为球队$i$所在地的常住人口（万人），$t_{ik}$为公路旅行时间（分钟）。该定义的含义是：将人口较多地区的球迷出行时间赋予更高权重，使选址倾向于服务人口密集区域。

\subsection{选址的等效成本扩展}

上述多维实验采用“替换成本矩阵或限制候选集”的方式扩展基准模型。若赛事组织方希望把更多软目标放入同一个目标函数，也可在本节基础上使用如下等效成本形式：

\begin{equation}
  c_{ik}^{\mathrm{ext}}
  = \alpha d_{ik} + \beta h_{ik} - \gamma I_k + \delta r_k + \eta b_{ik}.
\end{equation}

其中$d_{ik}$可取距离或旅行时间，$h_{ik}$为主场惩罚，$I_k$为城市影响力得分，$r_k$为区域失衡惩罚，$b_{ik}$为山区、海岛或跨湾交通不便修正项。本文当前实证只显式使用距离/时间成本、人口加权成本和影响力Top20候选集；其余项作为赛事组织细化时的扩展接口。具体含义如下：

\begin{enumerate}[(1)]
  \item \textbf{主场公平}：本文已将主场禁止作为硬约束；若实际组织中允许有限主场因素，也可改为惩罚项$h_{ik}$。
  \item \textbf{城市影响力}：优先选择地级市和经济强县作为场地（$I_k$项），或直接限制候选集为人口与收入综合得分最高的Top20地点，以扩大赛事传播效果。
  \item \textbf{区域覆盖}：通过$r_k$项惩罚区域集中，确保浙北、浙中、浙南和沿海地区均有场地分布。
\end{enumerate}

\section{问题4：赛制合理化建议}

\subsection{仿真模型}

\subsubsection{Bradley-Terry模型}

设球队$i$的实力为$s_i$，则球队$i$战胜球队$j$的概率为：

\begin{equation}
  P(i \text{ 胜 } j) = \frac{s_i}{s_i + s_j}
  \label{eq:bt}
\end{equation}

\subsubsection{同分排序与进球数扰动}

主仿真聚焦赛制公平性，因此单场只判胜负，不引入平局。真实足球小组赛采用积分3/1/0制，平局是常见的比赛结果；当前“必分胜负”的简化会高估单场比赛的排序信息量。本文将含平局积分制作为后续扩展，正文只比较同一简化模型下单循环和双循环的相对差异。小组赛积分相同时，用进球数作为第二排序指标；为避免完全确定性的并列排序，令球队$i$对阵球队$j$时的进球贡献近似服从Poisson分布：

\begin{equation}
  G_{ij} \sim \text{Poisson}(\lambda_{ij}),\quad \lambda_{ij} = 0.5 + 0.2\,s_i - 0.1\,s_j.
  \label{eq:poisson}
\end{equation}

进球期望$\lambda_{ij}$同时取决于攻方实力$s_i$和守方实力$s_j$：强队对弱队时$\lambda$更高，强队对强队时$\lambda$更低，符合足球比赛的攻守对抗规律。在本题所有实力取值下$\lambda_{ij} \geq 0.1$，保证Poisson参数始终为正。该Poisson项只用于同分时的进球数比较，不改变Bradley-Terry模型给出的胜负概率。

\subsubsection{实力赋值}

前文采用行政层级的三档固定赋值来估算实力，但固定三档赋值会导致同档球队完全并列，因此在这里的仿真采用层内随机化实力模型。每次完整赛事仿真前，按球队层级抽取潜在真实实力：

\begin{equation}
  s_i=\max\{\varepsilon,\ \mu_{\ell(i)}+\xi_i\},\quad
  \xi_i\sim N(0,\sigma^2),\quad \sigma=\frac{1}{3},
  \label{eq:random_strength}
\end{equation}

其中$\mu_{\ell(i)}\in\{3,2,1\}$分别对应市级、县级市和县队，$\varepsilon=0.05$用于保证Bradley-Terry概率中的实力值为正。相邻层级均值差为$3\sigma$，既保留层级强弱差异，也允许同档球队之间存在合理波动，使每次仿真中各队实力各不相同。

\subsection{评价指标}

\begin{enumerate}[(1)]
  \item \textbf{Spearman相关系数}：最终排名与真实实力排名的相关程度，衡量赛制对实力排序的保真度。由于淘汰赛只产生部分顺序，本文按比赛可观察信息补全名次：冠军列第1名，随后按淘汰轮次由后到前排列同轮出局球队；同轮出局球队不再用真实实力排序，而按当轮随机对阵顺序补位。小组赛未出线的32队按小组第3名、第4名分层列为33--64名，同层随机补位。不同补全规则会影响Spearman系数的绝对值，因此这里只在同一补位规则下比较赛制差异。
  \item \textbf{晋级准确率}：真实实力Top-32球队中实际晋级32强的比例。
  \item \textbf{冠军集中度}：统计实力Top-1、Top-4、Top-8球队的夺冠概率。
\end{enumerate}

\subsection{仿真结果}

对当前赛制（单循环小组赛+单场淘汰赛）和替代赛制（双循环小组赛+单场淘汰赛）各进行5\,000次完整赛事仿真，结果如表\ref{tab:q4_random_strength}。

\begin{table}[H]
  \centering
  \caption{赛制方案对比（层内正态随机实力模型，$\sigma=1/3$，5\,000次模拟）}
  \label{tab:q4_random_strength}
  \begin{tabular}{lcc}
    \toprule
    指标 & 单循环+淘汰（当前） & 双循环+淘汰（替代） \\
    \midrule
    Spearman相关系数（均值） & 0.471 & \textbf{0.564} \\
    Top-32晋级准确率 & 69.15\% & \textbf{74.59\%} \\
    Top-1夺冠率 & 7.58\% & 7.66\% \\
    Top-4夺冠率 & 26.76\% & 27.44\% \\
    Top-8夺冠率 & 47.12\% & 48.74\% \\
    Spearman 5\%分位 & 0.293 & 0.409 \\
    Spearman 95\%分位 & 0.630 & 0.702 \\
    每组比赛场次 & 6 & 12 \\
    总比赛场次 & $16\times6 + 31 = 127$ & $16\times12 + 31 = 223$ \\
    \bottomrule
  \end{tabular}
\end{table}

双循环显著改善小组出线和整体排名保真度，但对冠军集中度的提升有限。其原因是冠军仍由32强单场淘汰产生，淘汰赛偶然性成为冠军归属的主要扰动来源。

\subsection{分析与建议}

\subsubsection{当前赛制的优势与不足}

\textbf{优势}：
\begin{enumerate}[(1)]
  \item 赛程紧凑（127场比赛），组织成本低。
  \item 淘汰赛单场制悬念大、观赏性高。
  \item 16组同时开赛，初期关注面广。
\end{enumerate}

\textbf{不足}：
\begin{enumerate}[(1)]
  \item Spearman相关系数仅0.471，最终名次与实力排序只有中等程度一致，赛制对真实实力的保真度仍有提升空间。
  \item Top-32晋级准确率仅约69.2\%，约三成实力前32强球队会因小组赛偶然性无缘淘汰赛。
  \item 单循环每队仅3场比赛，样本量小导致小组排名偶然性大。
\end{enumerate}

\subsubsection{改进建议}

表\ref{tab:q4_random_strength}只对”小组赛单循环/双循环”这一核心变量进行了可复现实证比较。以下建议中，”小组赛改为双循环”有直接仿真支撑；其余建议属于基于仿真诊断和赛事组织逻辑提出的结构性增强方案，可在后续版本中继续通过完整赛程仿真验证。

{\renewcommand{\labelenumi}{\textbf{建议\arabic{enumi}：}}
  \begin{enumerate}

    \item \textbf{小组赛改为双循环}：Spearman相关系数从0.471提升至0.564，Top-32晋级准确率从69.2\%提升至74.6\%。代价是比赛场次增加96场（约75.6\%），需权衡赛事周期和成本。

    \item \textbf{引入种子分档抽签}：在问题2抽签流程的基础上增加分档约束。可按历史成绩、行政层级和上届成绩将64支队伍分为4档，每档16支，抽签时每组各取一档球队，同时继续保留C1、C2、C3地域约束。实现上可将“每组每档1队”加入前瞻流的剩余容量状态，先筛除会破坏分档容量或地域约束的候选组，再进行现场抽取。该机制有助于减少极端“死亡之组”，提高小组赛和淘汰赛阶段的竞技均衡。

    \item \textbf{采用同组回避的交叉对阵}：淘汰赛首轮采用“小组第1名 vs 其他组第2名”的交叉规则，并保证同组球队首轮不相遇；若采用种子排名，还可避免综合实力排名靠前的队伍过早相遇。这能在不显著增加场次的情况下提升淘汰赛后半程的比赛质量。

    \item \textbf{淘汰赛关键轮次采用两回合制}：随机实力模型中，双循环对Top-8夺冠率的影响有限（表\ref{tab:q4_random_strength}），说明冠军偶然性主要来自单场淘汰。从八强或四强阶段引入主客场两回合制可进一步降低关键轮次偶然性。该建议属于赛制结构外推，完整的定量验证需在仿真中显式加入两回合总比分、加时或点球规则，本文不直接给出两回合晋级概率。

    \item \textbf{增设小组第三附加赛与排位赛}：16个小组第三名可进行单败附加赛，决出前4名并设置“优胜杯”或专项奖励，但不影响冠军争夺。该设计能让中游球队在小组赛后仍有明确目标，并提高县级队和基层球迷的参与感。小组第四名及未进入附加赛的队伍也可安排短程排位赛或交流赛，扩大赛事覆盖面。

    \item \textbf{双败淘汰作为试点增强方案}：若赛事周期和组织能力允许，可在32强阶段试点双败淘汰赛。胜者组负者进入败者组，败者组优胜者与胜者组冠军会师决赛；若胜者组冠军在总决赛失利，可设置一次加赛。双败机制能减少“一场失误即出局”的偶然性，但赛程、转播和观众理解成本均高，建议先作为试点或在八强以后局部采用。

    \item \textbf{动态平衡赛程}：根据问题3得到的旅行距离和旅行时间矩阵，合理安排比赛间隔，避免某队短期内连续长途旅行。双循环、附加赛或双败赛制若被采用，应同步做赛程整数规划，使场地利用、休息时间和出行负担保持均衡。

\end{enumerate}}

\subsubsection{分阶段实施路径}

综合公平性收益和组织成本，建议采用分阶段落地方案。第一阶段优先实施双循环小组赛，这是当前仿真直接支持、且规则解释成本较低的改动；同组回避交叉对阵可作为低成本配套规则同步试行。第二阶段引入种子分档抽签和小组第三附加赛，用较小的额外组织成本改善小组均衡和参赛体验。第三阶段再试点双败淘汰或关键轮次两回合制，用于进一步提高冠军归属与真实实力的一致性。

\section{结论}

本文围绕浙江省市县联赛的赛事组织问题，建立了从分组、抽签、选址到赛制评价的完整数学模型与算法框架。主要贡献如下：

\begin{enumerate}[(1)]
  \item \textbf{分组方案}：建立了CSP模型，提出贪心、蛇形轮转和ILP四种分组策略。ILP均衡方案实现了$F_1=0$、$F_2=0.484$、$F_2'=1$的最优均衡性能。

  \item \textbf{抽签方案}：设计了前瞻最小费用流辅助抽签法，将C1、C2硬约束与C3最小冲突目标统一到网络流前向校验中。该方案在每次抽出县级队后只开放能够延续全局最优剩余分配的小组集合，既避免现场反悔换组，也在配对实验中实现零失败和$F_1=0$。

  \item \textbf{比赛地点}：建立了带主场禁止约束的设施选址+容量指派ILP模型，在haversine距离成本下得到总旅行距离7\,801\,km的最优方案。多维成本实验表明，交通时间、人口加权和影响力候选集会改变具体场地集合，但诸暨、杭州、义乌、绍兴等高频场地具有较强稳健性。

  \item \textbf{赛制建议}：通过5\,000次赛事仿真，定量证明双循环小组赛可将Spearman相关系数从0.471提升至0.564，Top-32晋级准确率从69.2\%提升至74.6\%。结合冠军集中度结果，本文进一步建议引入种子分档、同组回避交叉对阵、小组第三附加赛、关键轮次两回合制和双败淘汰试点。

  \item \textbf{抽签鲁棒性验证}：在问题2中通过100个随机行政拓扑比较基础贪心与前瞻网络流，发现前瞻流在全部配对实验中均未死锁；当C3零冲突不可实现时，网络流模型会自动转为最小冲突分配，进一步支持采用前瞻流进行事前筛选。
\end{enumerate}

本文建立的模型框架可以为类似层级结构体育赛事的分组与赛制设计提供方法论参考。

\newpage

\section*{附录A：分组方案详细结果}

\subsection*{方案A（贪心优先填充）}

{\small
\begin{verbatim}
组 1 [实力=7] 温州市(市), 遂昌县(县), 浦江县(县), 余姚市(县市级)
组 2 [实力=8] 丽水市(市), 平阳县(县), 东阳市(县市级), 江山市(县市级)
组 3 [实力=7] 金华市(市), 永嘉县(县), 玉环市(县市级), 常山县(县)
组 4 [实力=7] 台州市(市), 文成县(县), 桐乡市(县市级), 龙游县(县)
组 5 [实力=6] 嘉兴市(市), 景宁县(县), 仙居县(县), 开化县(县)
组 6 [实力=8] 宁波市(市), 龙泉市(县市级), 天台县(县), 建德市(县市级)
组 7 [实力=7] 衢州市(市), 松阳县(县), 温岭市(县市级), 桐庐县(县)
组 8 [实力=6] 杭州市(市), 云和县(县), 三门县(县), 安吉县(县)
组 9 [实力=7] 湖州市(市), 缙云县(县), 临海市(县市级), 淳安县(县)
组10 [实力=6] 绍兴市(市), 青田县(县), 嘉善县(县), 德清县(县)
组11 [实力=7] 舟山市(市), 庆元县(县), 海宁市(县市级), 长兴县(县)
组12 [实力=6] 乐清市(县市级), 兰溪市(县市级), 海盐县(县), 新昌县(县)
组13 [实力=6] 苍南县(县), 磐安县(县), 平湖市(县市级), 诸暨市(县市级)
组14 [实力=8] 瑞安市(县市级), 永康市(县市级), 慈溪市(县市级), 嵊州市(县市级)
组15 [实力=5] 泰顺县(县), 义乌市(县市级), 象山县(县), 岱山县(县)
组16 [实力=5] 龙港市(县市级), 武义县(县), 宁海县(县), 嵊泗县(县)

指标: F1=0  F2=0.927  F2'=3
\end{verbatim}
}

\subsection*{方案B（蛇形轮转分配）}

{\small
\begin{verbatim}
组 1 [实力=8] 温州市(市), 东阳市(县市级), 玉环市(县市级), 长兴县(县)
组 2 [实力=7] 丽水市(市), 乐清市(县市级), 浦江县(县), 仙居县(县)
组 3 [实力=6] 金华市(市), 苍南县(县), 天台县(县), 德清县(县)
组 4 [实力=7] 台州市(市), 瑞安市(县市级), 武义县(县), 安吉县(县)
组 5 [实力=8] 嘉兴市(市), 泰顺县(县), 义乌市(县市级), 温岭市(县市级)
组 6 [实力=8] 宁波市(市), 龙港市(县市级), 永康市(县市级), 三门县(县)
组 7 [实力=7] 衢州市(市), 平阳县(县), 磐安县(县), 临海市(县市级)
组 8 [实力=8] 杭州市(市), 永嘉县(县), 兰溪市(县市级), 桐乡市(县市级)
组 9 [实力=6] 湖州市(市), 文成县(县), 庆元县(县), 嘉善县(县)
组10 [实力=7] 绍兴市(市), 遂昌县(县), 海宁市(县市级), 淳安县(县)
组11 [实力=6] 舟山市(市), 景宁县(县), 海盐县(县), 桐庐县(县)
组12 [实力=7] 龙泉市(县市级), 平湖市(县市级), 建德市(县市级), 新昌县(县)
组13 [实力=6] 松阳县(县), 慈溪市(县市级), 开化县(县), 诸暨市(县市级)
组14 [实力=5] 云和县(县), 象山县(县), 龙游县(县), 嵊州市(县市级)
组15 [实力=4] 缙云县(县), 宁海县(县), 常山县(县), 岱山县(县)
组16 [实力=6] 青田县(县), 余姚市(县市级), 江山市(县市级), 嵊泗县(县)

指标: F1=0  F2=1.111  F2'=4
\end{verbatim}
}

\subsection*{方案C（ILP-C3最小化）}

{\small
\begin{verbatim}
组 1 [实力=8] 宁波市(市), 平湖市(县市级), 东阳市(县市级), 天台县(县)
组 2 [实力=6] 象山县(县), 泰顺县(县), 金华市(市), 仙居县(县)
组 3 [实力=5] 永嘉县(县), 嵊州市(县市级), 龙游县(县), 松阳县(县)
组 4 [实力=6] 平阳县(县), 长兴县(县), 台州市(市), 庆元县(县)
组 5 [实力=5] 慈溪市(县市级), 嘉善县(县), 新昌县(县), 嵊泗县(县)
组 6 [实力=6] 宁海县(县), 桐乡市(县市级), 德清县(县), 江山市(县市级)
组 7 [实力=6] 桐庐县(县), 苍南县(县), 磐安县(县), 丽水市(市)
组 8 [实力=8] 建德市(县市级), 舟山市(市), 临海市(县市级), 遂昌县(县)
组 9 [实力=8] 瑞安市(县市级), 嘉兴市(市), 义乌市(县市级), 开化县(县)
组10 [实力=6] 乐清市(县市级), 安吉县(县), 常山县(县), 玉环市(县市级)
组11 [实力=8] 龙港市(县市级), 湖州市(市), 永康市(县市级), 云和县(县)
组12 [实力=6] 淳安县(县), 温州市(市), 武义县(县), 景宁县(县)
组13 [实力=7] 文成县(县), 诸暨市(县市级), 兰溪市(县市级), 龙泉市(县市级)
组14 [实力=7] 余姚市(县市级), 海盐县(县), 衢州市(市), 青田县(县)
组15 [实力=7] 海宁市(县市级), 绍兴市(市), 三门县(县), 缙云县(县)
组16 [实力=7] 杭州市(市), 浦江县(县), 岱山县(县), 温岭市(县市级)

指标: F1=0  F2=0.992  F2'=3
\end{verbatim}
}

\subsection*{方案D（ILP均衡优化）—— 推荐方案}

{\small
\begin{verbatim}
组 1 [实力=6] 龙港市(县市级), 德清县(县), 东阳市(县市级), 三门县(县)
组 2 [实力=6] 嘉善县(县), 金华市(市), 岱山县(县), 庆元县(县)
组 3 [实力=7] 永嘉县(县), 平湖市(县市级), 长兴县(县), 舟山市(市)
组 4 [实力=7] 温州市(市), 海盐县(县), 诸暨市(县市级), 磐安县(县)
组 5 [实力=7] 余姚市(县市级), 安吉县(县), 浦江县(县), 丽水市(市)
组 6 [实力=7] 杭州市(市), 象山县(县), 兰溪市(县市级), 常山县(县)
组 7 [实力=7] 宁波市(市), 江山市(县市级), 仙居县(县), 青田县(县)
组 8 [实力=6] 建德市(县市级), 苍南县(县), 龙游县(县), 温岭市(县市级)
组 9 [实力=7] 武义县(县), 衢州市(市), 玉环市(县市级), 云和县(县)
组10 [实力=7] 文成县(县), 绍兴市(市), 永康市(县市级), 开化县(县)
组11 [实力=7] 宁海县(县), 嘉兴市(市), 义乌市(县市级), 遂昌县(县)
组12 [实力=6] 瑞安市(县市级), 海宁市(县市级), 新昌县(县), 松阳县(县)
组13 [实力=6] 桐庐县(县), 乐清市(县市级), 临海市(县市级), 景宁县(县)
组14 [实力=7] 淳安县(县), 慈溪市(县市级), 台州市(市), 缙云县(县)
组15 [实力=7] 平阳县(县), 桐乡市(县市级), 嵊州市(县市级), 龙泉市(县市级)
组16 [实力=6] 泰顺县(县), 湖州市(市), 嵊泗县(县), 天台县(县)

指标: F1=0  F2=0.484  F2'=1
\end{verbatim}
}

\subsection*{ILP最优选址方案（基于方案D分组）}

{\small
\begin{verbatim}
场地1: 余姚市  承办 组2、组3    场地2: 新昌县  承办 组12、组16
场地3: 东阳市  承办 组1、组11   场地4: 永康市  承办 组10、组15
场地5: 武义县  承办 组8、组9    场地6: 浦江县  承办 组5、组6
场地7: 磐安县  承办 组4、组14   场地8: 仙居县  承办 组7、组13

总旅行距离: 7801 km | 平均: 121.9 km/队 | 最远: 337 km
\end{verbatim}
}

\section*{附录B：代码清单}
\label{app:code}

所有代码使用Python 3实现，依赖NumPy、SciPy和PuLP（调用COIN-OR CBC求解器）。

\subsection*{B.1 问题1：分组方案}

\lstinputlisting[caption={q1\_grouping.py\ ——\ 球队数据、约束检查、四种分组方案（贪心、蛇形轮转、ILP-C3最小化、ILP均衡优化）及$F_1,F_2,F_2'$三项评价指标}]{q1_grouping.py}

\subsection*{B.2 问题2：前瞻抽签与鲁棒性检验}

\lstinputlisting[caption={q2\_lookahead\_flow.py\ ——\ Dinic最大流与最小费用最大流实现，前瞻辅助抽签法，含C3零冲突快速最大流检查}]{q2_lookahead_flow.py}

\lstinputlisting[caption={strategy\_comparison.py\ ——\ 朴素贪心与前瞻流在1\,000个配对种子上的策略对比实验，含边际均匀性卡方检验}]{strategy_comparison.py}

\lstinputlisting[caption={generalized\_mc.py\ ——\ 随机行政拓扑生成与基础贪心、前瞻网络流的广义蒙特卡洛鲁棒性验证}]{generalized_mc.py}


\subsection*{B.3 问题3：比赛地点选择}

\lstinputlisting[caption={q3\_venue.py\ ——\ Haversine距离、公路/铁路旅行时间、人口加权等多维成本模型的ILP设施选址，含主场禁止约束和启发式对比}]{q3_venue.py}

\subsection*{B.4 问题4：赛制仿真}

\lstinputlisting[caption={q4\_simulation.py\ ——\ Bradley-Terry胜负模型与层内正态随机实力下的5\,000次赛事蒙特卡洛仿真}]{q4_simulation.py}

\subsection*{B.5 实验输出文件}

正文各表格取数来源如下：

\begin{table}[H]
  \centering
  \caption*{输出文件与对应表格}
  \small
  \setlength{\tabcolsep}{3pt}
  \begin{tabular}{p{6.5cm}p{8.5cm}}
    \toprule
    输出文件 & 对应表格 \\
    \midrule
    \texttt{q1\_output.txt} & 表\ref{tab:q1_comparison}（四种分组方案指标对比） \\
    \texttt{strategy\_comparison\_output.txt} & 表\ref{tab:strategy_comparison}（朴素贪心与前瞻流1\,000种子实验） \\
    \texttt{generalized\_mc\_output.txt} & 表\ref{tab:generalized}（100个随机拓扑的广义蒙特卡洛） \\
    \texttt{q3\_output.txt} & 表\ref{tab:q3_ilp}（Haversine距离选址结果） \\
    \texttt{q3\_output\_v2.txt} & 多维交通成本选址实验输出 \\
    \texttt{q4\_output.txt} & 表\ref{tab:q4_random_strength}（5\,000次赛制仿真结果） \\
    \bottomrule
  \end{tabular}
\end{table}

\end{document}
